\section{The Clock And The Trial}

Time enters the grammar as alternation. The foundation does not need a
cosmic container in which events float. It needs a readable flip: before
and after, tick and tock, open and closed, unresolved and resolved. The
alternation is binary because it must first be distinguishable. A clock
that cannot tell its two faces apart is not a clock. A clock with more
faces may be useful later, but the first temporal act is the permitted
return from one side of a difference to the other.

This makes time a grammatical achievement rather than an imported
substance. The grammar permits a reading to be placed before another
reading. It forbids the reader from treating all readings as
simultaneously available. It carries a local memory: enough of the
previous state to say that the next state has occurred. Once that memory
is available, a finite condition can become a trial.

A trial is not a single reading. It is a reading that can answer the
question, ``again?'' The question matters because measurement differs
from mere impression by repeatability. A flash in the dark may be
noticed, but until the grammar can specify how the noticing could be
returned to, the flash remains an event rather than a measurement. A
trial binds stimulus, expectation, and response into a repeatable form.
It says what may count as the same trial and what must count as a new or
failed trial.

The binary clock is therefore not a decorative addition to the count. It
is the condition under which count becomes experimental. A whole-reading
can now be repeated across time: the same stimulus, the same expected
kind of response, the same tolerance, the same carried unit. A
part-reading can now be located in time: the response changed here, the
tolerance failed there, the residue turned after this many steps. The
clock permits both preservation and fracture to become observable across
successive trials.

The Time Effect gives a useful illustration. An open record contains an
unsettled degree of freedom. When the record closes, an event has become
admissible. Time is read as the count of such closures, not as a hidden
fluid that pushed them along. The grammar need not assert that all
physical time is nothing but bookkeeping. It only insists that measured
time, the time that can enter the ledger, appears when an open reading is
resolved and a new reading can be placed after it.

This is why repeatability is stronger than recurrence. Recurrence says
that something has happened again. Repeatability says that the grammar
knows what would make it the same kind of happening. The difference is
crucial. Without repeatability, a count of events is only a list. With
repeatability, the list has a rule of return. The grammar can ask
whether the later reading preserves the earlier whole or exposes a new
part. It can compare trials because it has bounded what a trial is
allowed to carry.

The clock also makes irreversibility visible. A reading placed after
another reading cannot be unread as though it never entered the ledger.
The grammar may revise, quotient, or reinterpret records, but the
sequence of commitment remains. Measurement is a historical act. A
trial's later response may correct an expectation, but it cannot be the
same as never having asked the question.

The trial therefore binds the first four obligations of the chapter. It
uses distinguishability to separate clock faces. It uses counting to
carry repetition. It uses finite conditions to avoid completed
omniscience. It uses residue to decide what remains after each
resolution. What it cannot do is spend unlimited time. Every trial
inherits a horizon.
